\documentclass[11pt]{article}
\usepackage[utf8]{inputenc}
\usepackage{amsmath, amssymb, amsthm}
\usepackage[margin=1in]{geometry}

\theoremstyle{definition}
\newtheorem{definition}{Definition}
\newtheorem{theorem}{Theorem}
\newtheorem{lemma}{Lemma}

\theoremstyle{remark}
\newtheorem{remark}{Remark}

\DeclareMathOperator{\tr}{tr}
\newcommand{\dual}[2]{{}_{V'}\langle #1, #2 \rangle_{V}}

\title{Lecture 1: Introduction to SPDEs and the Heat Equation}
\date{16.10.2025}
\author{Tien Anh Vu}

\begin{document}
\maketitle

Good morning. Let us begin our introduction to Stochastic Partial Differential Equations, or SPDEs.

Today, we will build some intuition by looking at a classic example that motivates the Dean-Kawasaki equation, originally introduced in the works of Dean in 1996 and Kawasaki in 1998. We will start from a microscopic particle system and observe how it leads us to a macroscopic partial differential equation.

\section*{1. Microscopic Particle System}
Let us consider a system of independent Brownian motions, which we will denote as $B^1, B^2, B^3, \ldots$. We also take a test function $f$ that is smooth. Specifically, we require $f\in C_b^2(\mathbb{R})$, which means that $f$ is twice continuously differentiable and has bounded derivatives.

\begin{definition}[Empirical average]
We can define the empirical average of this system of $N$ particles at time $t$ as $\eta_t^N(f)$. We write this as:
\[
\eta_t^N(f):=\frac{1}{N}\sum_{i=1}^{N} f\bigl(B_t^i+x_0\bigr).
\]
Different choices of $f$ probe different features:
\begin{itemize}
\item if $f(x)=x$, then we are averaging particle positions (mean position);
\item if $f(x)=x^2$, then we are averaging squared positions (second moment/spread);
\end{itemize}
If we want to emphasize sample-wise evaluation, we write
\[
\eta_t^N(f)(\omega):=\frac{1}{N}\sum_{i=1}^{N} f\bigl(B_t^i(\omega)+x_0\bigr).
\]
where $x_0$ represents the starting position of our particles. We can equivalently express this empirical average as an integral with respect to the empirical measure, utilizing the Dirac delta measure $\delta$:
\[
\eta_t^N(f)=\int f\,d\!\left(\frac{1}{N}\sum_{i=1}^{N}\delta_{B_t^i+x_0}\right).
\]
This is shorthand for the same integral with the variable shown explicitly:
\[
\eta_t^N(f)=\int_{\mathbb R} f(x)\,d\!\left(\frac{1}{N}\sum_{i=1}^{N}\delta_{B_t^i+x_0}\right)(x).
\]
\end{definition}

\section*{2. Macroscopic Limit via SLLN}
Now, we want to understand the macroscopic behavior of this system as the number of particles grows very large. We apply the Strong Law of Large Numbers (SLLN). As $N\to\infty$, this empirical average converges almost surely with respect to $P$-measure to the expected value:
\[
\eta_t^N(f)\xrightarrow[N\to\infty]{\text{a.s. w.r.t. }P}\mathbb{E}\bigl[f(B_t+x_0)\bigr].
\]

Since $B_t$ is a standard Brownian motion, a clean way to write this is
\[
B_t+x_0\sim \mathcal N(x_0,t),\qquad \mu_t(dx)=p(t,x)\,dx,
\]
with
\[
p(t,x)=\frac{1}{\sqrt{2\pi t}}\exp\!\left(-\frac{(x-x_0)^2}{2t}\right).
\]
Here $\mu_t$ is the probability distribution of $B_t+x_0$ at time $t$, and $p(t,x)$ is the density of $\mu_t$ (equivalently, of $B_t+x_0$) with respect to Lebesgue measure.
Then the expectation becomes
\[
\mathbb{E}\bigl[f(B_t+x_0)\bigr]=\int_{\mathbb{R}} f(x)p(t,x)\,dx.
\]
Intuitively, we are averaging the values of $f$ over all positions $x$, weighted by how likely the Brownian particle is to be at $x$ at time $t$.

\section*{3. Governing PDE: Heat Equation}
The natural question now is to identify the governing equation for this probability density $p(t,x)$. It turns out that $p(t,x)$ is the fundamental solution of the Heat Equation. The system is:
\[
\partial_t p(t,x)=\frac{1}{2}\partial_{xx}p(t,x),\qquad t>0,\;x\in\mathbb{R},
\]
coupled with the initial condition as $t\to0$:
\[
p(t,x)dx\to\delta_{x_0}(dx).
\]

\section*{4. Physical Intuition}
To build a physical intuition for the heat equation, consider the relationship between the time derivative $\partial_t p$ and the spatial second derivative $\partial_{xx}p$. If you plot the density $p(t,x)$ and look at a region with negative curvature---such as a local peak---the second spatial derivative $\partial_{xx}p(t,x)$ is negative. According to the heat equation, this means the time derivative $\partial_t p(t,x)$ must also be negative. Consequently, at points of negative curvature, the ``temperature'' or probability density will go down, which has the effect of smoothing out the distribution over time.

We will stop here for this step of the derivation. Are there any questions on the transition from the empirical average to the heat equation?

\section*{5. Microscopic Fluctuations and the Macroscopic Limit}
\subsection*{5.1. Setup and Empirical Average}
We have already introduced the microscopic particle system, defined the empirical average, and identified the corresponding macroscopic heat equation. We now continue by making that passage more rigorous: we analyze the fluctuations around the Law of Large Numbers using It\^o's formula and show how the stochastic part disappears as the number of particles grows.

To do this, we return to the empirical measure for $N$ independent Brownian motions $B_t^i$, all starting at an initial position $x_0$, and evaluate it against a test function $f$. This gives the empirical average
\[
\eta_t^N(f)=\frac{1}{N}\sum_{i=1}^N f(B_t^i+x_0).
\]

\subsection*{5.2. It\^o Formula and the Drift-Martingale Decomposition}
To observe the time evolution of this empirical measure, we compute its stochastic differential. Applying It\^o's formula, the differential is expressed as:
\[
d\!\left(
\frac{1}{N}\sum_{i=1}^N f(B_t^i+x_0)
\right)
=
\frac{1}{N}\sum_{i=1}^N df(B_t^i+x_0).
\]

When we expand $df(B_t^i+x_0)$ using the standard rules of stochastic calculus, we decompose the dynamic into two fundamentally distinct components: an It\^o integral and a correction term (the drift). The expansion gives us:
\[
=
\frac{1}{N}\sum_{i=1}^N f'(B_t^i+x_0)\,dB_t^i
+
\frac{1}{N}\sum_{i=1}^N \frac{1}{2}f''(B_t^i+x_0)\,dt.
\]

Let us analyze these two terms closely. The second term is our correction term, which can be written compactly using our empirical measure notation as $\eta_t^N\!\left(\frac{1}{2}f''\right)$. This term captures the expected diffusion, effectively mapping to the Laplacian operator in our emerging macroscopic PDE.

The first term is a local Martingale, which we will denote as $M_t^N(f)$. This term represents the microscopic stochastic fluctuations of the system. To understand the macroscopic limit, we must determine the asymptotic behavior of this It\^o integral as $N\to\infty$. Specifically, we need to prove that
\[
M_t^N(f)\xrightarrow[N\to\infty]{}0.
\]

\subsection*{5.3. Vanishing of the Martingale via It\^o Isometry}
To prove that the fluctuations vanish, we compute the expected value of the squared martingale, $\mathbb E[(M_t^N(f))^2]$. We do this by applying It\^o isometry, which allows us to equate the expectation of a squared stochastic integral to the expectation of the standard integral of the squared integrand. Because the Brownian motions are independent, the cross-terms in the squared sum take an expectation of zero, leaving us with:
\[
\mathbb E[(M_t^N(f))^2]
=
\frac{1}{N^2}\sum_{i=1}^N
\mathbb E\!\left[
\left(\int_0^t f'(B_s^i+x_0)\,dB_s^i\right)^2
\right].
\]

By It\^o isometry, we transform the stochastic integral into a Riemann integral:
\[
=
\frac{1}{N^2}\sum_{i=1}^N
\mathbb E\!\left[
\int_0^t f'(B_s^i+x_0)^2\,ds
\right].
\]

Since the $N$ Brownian motions are identically distributed, the expectation is the same for every particle $i$. This allows us to simplify the sum completely:
\[
=
\frac{1}{N}
\mathbb E\!\left[
\int_0^t f'(B_s+x_0)^2\,ds
\right].
\]

\subsection*{5.4. Quadratic Variation}
We can apply the exact same logic to compute the quadratic variation of the martingale, $\langle M^N(f)\rangle_t$. Evaluating the quadratic variation yields:
\[
\langle M^N(f)\rangle_t
=
\frac{1}{N^2}\sum_{i=1}^N
\int_0^t f'(B_s^i+x_0)^2\,ds.
\]

This expression can be rewritten by pulling the $\frac{1}{N}$ factor out and recognizing the remaining sum as the empirical measure applied to $(f')^2$. Thus, we obtain:
\[
=
\frac{1}{N}\int_0^t \eta_s^N((f')^2)\,ds.
\]

The critical insight here is that the quadratic variation is of order $O\!\left(\frac{1}{N}\right)$. Because the variance of the martingale term scales by $\frac{1}{N}$, it completely vanishes in the limit as the number of molecules $N$ approaches infinity.

\subsection*{5.5. Conclusion}
This demonstrates the Law of Large Numbers for our empirical measure. The stochastic noise that characterizes the microscopic gas molecules dissipates entirely, leaving only the deterministic correction term. This remaining term dictates the macroscopic description of the system, effectively deriving the governing PDE from the underlying stochastic particle dynamics.

\section*{Appendix}
\subsection*{A.1. Brownian Motion}
\begin{definition}[Brownian motion]
Let $(\Omega,\mathcal F,(\mathcal F_t)_{t\ge0},\mathbb P)$ be a filtered probability space. A process $(W_t)_{t\ge0}$ is a standard Brownian motion if:
\begin{itemize}
\item $W_0=0$ almost surely;
\item for every $0\le s<t$, $W_t-W_s\sim\mathcal N(0,t-s)$;
\item increments on disjoint time intervals are independent;
\item sample paths are continuous.
\end{itemize}
In this lecture, the particle position is modeled by $X_t^i=B_t^i+x_0$, where each $B^i$ is a Brownian motion.
\end{definition}

\subsection*{A.2. Empirical Average}
\begin{definition}[General empirical average]
Let $E$ be a measurable space, let $Y^1,\dots,Y^N$ be $E$-valued random variables, and let $f:E\to\mathbb R$ be bounded measurable. The empirical average is the random variable
\[
\eta^N(f):=\frac{1}{N}\sum_{i=1}^N f(Y^i).
\]
For a fixed realization $\omega$, this becomes the realized empirical average
\[
\eta^N(f)(\omega):=\frac{1}{N}\sum_{i=1}^N f\bigl(Y^i(\omega)\bigr).
\]
This quantity is the sample mean of the observable $f$ evaluated on the $N$ particles. Thus, when we simply say ``empirical average,'' we usually mean the random variable $\eta^N(f)$, while $\eta^N(f)(\omega)$ denotes the realized average for one sample $\omega$.
In this lecture, we take
\[
Y^i=X_t^i:=B_t^i+x_0,
\]
so
\[
\eta_t^N(f)=\frac{1}{N}\sum_{i=1}^N f(B_t^i+x_0),
\]
and for a fixed sample $\omega$,
\[
\eta_t^N(f)(\omega)=\frac{1}{N}\sum_{i=1}^N f\bigl(B_t^i+x_0\bigr)(\omega).
\]
Thus $\eta_t^N(f)$ is the particle-system empirical average of $f$ at time $t$.
\end{definition}

\subsection*{A.3. Empirical Measure}
\begin{definition}[Empirical measure]
For $Y^1,\dots,Y^N$ as above, the empirical measure is the random measure
\[
\mu^N:=\frac{1}{N}\sum_{i=1}^N \delta_{Y^i}.
\]
For a fixed realization $\omega$, this becomes the realized empirical measure
\[
\mu^N(\omega):=\frac{1}{N}\sum_{i=1}^N \delta_{Y^i(\omega)}.
\]
Then
\[
\eta^N(f)(\omega)=\int_E f(y)\,\mu^N(\omega)(dy).
\]
Thus, the empirical average is obtained by testing the empirical measure against the function $f$. In this lecture, at time $t$,
\[
\mu_t^N:=\frac{1}{N}\sum_{i=1}^N \delta_{B_t^i+x_0},
\]
and for a fixed sample $\omega$,
\[
\mu_t^N(\omega):=\frac{1}{N}\sum_{i=1}^N \delta_{B_t^i(\omega)+x_0}.
\]
Then
\[
\eta_t^N(f)(\omega)=\int_{\mathbb R} f(x)\,\mu_t^N(\omega)(dx).
\]
\end{definition}

\subsection*{A.4. Strong Law of Large Numbers (SLLN)}
\begin{theorem}[SLLN]
Let $(Z_i)_{i\ge1}$ be i.i.d.\ random variables with $\mathbb E[|Z_1|]<\infty$. Then
\[
\frac{1}{N}\sum_{i=1}^N Z_i
\xrightarrow[N\to\infty]{\text{a.s.}}
\mathbb E[Z_1].
\]
This says that the sample average converges almost surely to the true expectation.
\end{theorem}

\begin{remark}[Application in this lecture]
For fixed $t$ and fixed test function $f$, take
\[
Z_i:=f(B_t^i+x_0).
\]
Since $(B_t^i)_i$ are independent and identically distributed, so are $(Z_i)_i$. Therefore
\[
\eta_t^N(f)
=
\frac{1}{N}\sum_{i=1}^N f(B_t^i+x_0)
\xrightarrow[N\to\infty]{\text{a.s.}}
\mathbb E[f(B_t+x_0)].
\]
This is exactly the convergence used in Section 2.
\end{remark}

\subsection*{A.5. Types of Convergence for Random Variables}
\begin{definition}[Almost sure convergence]
We say $X_n\to X$ almost surely (a.s.) if
\[
\mathbb P\Bigl(\bigl\{\omega:\,\lim_{n\to\infty}X_n(\omega)=X(\omega)\bigr\}\Bigr)=1.
\]
In words: for almost every outcome $\omega$, if we look at the sequence $X_n(\omega)$, it converges pointwise to $X(\omega)$.
\end{definition}

\begin{definition}[Convergence in probability]
We say $X_n\to X$ in probability if for every $\varepsilon>0$,
\[
\lim_{n\to\infty}\mathbb P(|X_n-X|>\varepsilon)=0.
\]
So if you fix any error tolerance $\varepsilon$, the chance that $X_n$ is farther than $\varepsilon$ from $X$ goes to zero as $n$ increases.
\end{definition}

\begin{definition}[Convergence in $L^p$]
For $p\ge1$, we say $X_n\to X$ in $L^p$ if
\[
\lim_{n\to\infty}\mathbb E\bigl[|X_n-X|^p\bigr]=0.
\]
You can think of this as mean convergence with a $p$-power penalty on the error; larger $p$ puts more weight on large deviations.
\end{definition}

\begin{definition}[Convergence in distribution]
We say $X_n\to X$ in distribution (or weakly) if
\[
\lim_{n\to\infty}\mathbb E[\phi(X_n)]=\mathbb E[\phi(X)]
\]
for every bounded continuous test function $\phi$.
Equivalently, if $F_n$ and $F$ are the distribution functions of $X_n$ and $X$, then
\[
\lim_{n\to\infty}F_n(x)=F(x)
\]
at every continuity point $x$ of $F$. Here we are not tracking sample paths directly; we are saying that the probability laws themselves converge.
\end{definition}

\begin{remark}[Standard implications]
The main implication chain is
\[
X_n\to X\text{ in }L^p
\;\Longrightarrow\;
X_n\to X\text{ in probability}
\;\Longrightarrow\;
X_n\to X\text{ in distribution}.
\]
Also, almost sure convergence implies convergence in probability.
When teaching this, a useful rule of thumb is: stronger notions (like a.s.\ or $L^p$) usually give weaker ones (probability, then distribution).
\end{remark}

\subsection*{A.6. Expectation}
\begin{definition}[Expectation]
Let $(\Omega,\mathcal F,\mathbb P)$ be a probability space and let $X:\Omega\to\mathbb R$ be an integrable random variable (that is, $\mathbb E[|X|]<\infty$). Its expectation is
\[
\mathbb E[X]:=\int_{\Omega} X(\omega)\,\mathbb P(d\omega).
\]
If $X$ has distribution $\mu_X$ on $\mathbb R$, then equivalently
\[
\mathbb E[X]=\int_{\mathbb R} x\,\mu_X(dx).
\]
More generally, for measurable $f$ with $\mathbb E[|f(X)|]<\infty$,
\[
\mathbb E[f(X)]=\int_{\mathbb R} f(x)\,\mu_X(dx).
\]
In our lecture notation, this is exactly the form
\[
\mathbb E[f(B_t+x_0)]=\int_{\mathbb R} f(x)\,p(t,x)\,dx.
\]
\end{definition}

\subsection*{A.7. Distribution $\mu_X$}
\begin{definition}[Distribution (probability law) of a random variable]
Let $X:(\Omega,\mathcal F,\mathbb P)\to\mathbb R$ be a random variable. Its distribution (or probability law, equivalently probability distribution) is the measure $\mu_X$ on $\mathbb R$ defined by
\[
\mu_X(A):=\mathbb P(X\in A),
\]
for every Borel set $A\subseteq\mathbb R$.

So ``distribution of $X$,'' ``probability law of $X$,'' ``probability distribution of $X$,'' and ``the probability measure induced by $X$'' all refer to this same measure $\mu_X$.
So $\mu_X$ tells us how the values of $X$ are distributed, without referring to the sample space directly.
\end{definition}

\begin{remark}
For any bounded measurable $f$,
\[
\int_{\mathbb R} f(x)\,\mu_X(dx)=\mathbb E[f(X)].
\]
This is why we can rewrite expectations as integrals with respect to $\mu_X$.
\end{remark}

\begin{remark}[Pushforward measure viewpoint]
The same measure can be written as
\[
\mu_X = \mathbb P\circ X^{-1},
\]
read as ``the pushforward of $\mathbb P$ by $X$.'' Concretely, this means
\[
(\mathbb P\circ X^{-1})(A)=\mathbb P\!\bigl(X^{-1}(A)\bigr)=\mathbb P(X\in A),
\]
which is exactly the definition of $\mu_X$ above.
\end{remark}

\subsection*{A.8. Density}
\begin{definition}[Density of a random variable]
Let $X$ be a real-valued random variable with law $\mu_X$. We say that $X$ has a density $p_X$ if
\[
\mu_X(A)=\int_A p_X(x)\,dx
\]
for every Borel set $A\subseteq\mathbb R$.
Equivalently, $\mu_X$ is absolutely continuous with respect to Lebesgue measure $\lambda$ (written $\mu_X\ll\lambda$), and $p_X$ is the Radon--Nikodym derivative:
\[
p_X=\frac{d\mu_X}{d\lambda},
\]
which describes how the probability mass is distributed over $\mathbb R$.
\end{definition}

\begin{remark}
If $X$ has density $p_X$, then for every bounded measurable $f$,
\[
\mathbb E[f(X)]=\int_{\mathbb R} f(x)\,p_X(x)\,dx.
\]
In our lecture, $X=B_t+x_0$ has Gaussian density
\[
p(t,x)=\frac{1}{\sqrt{2\pi t}}\exp\!\left(-\frac{(x-x_0)^2}{2t}\right),
\]
so
\[
\mathbb E[f(B_t+x_0)]=\int_{\mathbb R} f(x)\,p(t,x)\,dx.
\]
\end{remark}

\subsection*{A.9. Meaning of $\delta_{x_0}(dx)$}
\begin{definition}[Dirac measure at $x_0$]
The symbol $\delta_{x_0}$ denotes the Dirac measure concentrated at the point $x_0$. It is defined by
\[
\delta_{x_0}(A)=\mathbf 1_A(x_0)
\]
for every Borel set $A\subseteq\mathbb R$. Equivalently,
\[
\delta_{x_0}(A)=
\begin{cases}
1, & \text{if } x_0\in A,\\
0, & \text{if } x_0\notin A.
\end{cases}
\]
\end{definition}

\begin{remark}
Writing $\delta_{x_0}(dx)$ emphasizes that this is a measure in the variable $x$. Its key property is
\[
\int_{\mathbb R} f(x)\,\delta_{x_0}(dx)=f(x_0).
\]
So in our note,
\[
p(t,x)\,dx\to\delta_{x_0}(dx)\quad (t\downarrow0)
\]
means the probability law with density $p(t,x)$ converges to a point mass at $x_0$.
\end{remark}

\end{document}
